%%=============================================================================
%% Methodologie
%%=============================================================================

\chapter{\IfLanguageName{dutch}{Methodologie}{Methodology}}%
\label{ch:methodologie}

%% TODO: In dit hoofstuk geef je een korte toelichting over hoe je te werk bent
%% gegaan. Verdeel je onderzoek in grote fasen, en licht in elke fase toe wat
%% de doelstelling was, welke deliverables daar uit gekomen zijn, en welke
%% onderzoeksmethoden je daarbij toegepast hebt. Verantwoord waarom je
%% op deze manier te werk gegaan bent.
%% 
%% Voorbeelden van zulke fasen zijn: literatuurstudie, opstellen van een
%% requirements-analyse, opstellen long-list (bij vergelijkende studie),
%% selectie van geschikte tools (bij vergelijkende studie, "short-list"),
%% opzetten testopstelling/PoC, uitvoeren testen en verzamelen
%% van resultaten, analyse van resultaten, ...
%%
%% !!!!! LET OP !!!!!
%%
%% Het is uitdrukkelijk NIET de bedoeling dat je het grootste deel van de corpus
%% van je bachelorproef in dit hoofstuk verwerkt! Dit hoofdstuk is eerder een
%% kort overzicht van je plan van aanpak.
%%
%% Maak voor elke fase (behalve het literatuuronderzoek) een NIEUW HOOFDSTUK aan
%% en geef het een gepaste titel.

Dit hoofdstuk beschrijft de systematische en empirische aanpak die is gehanteerd om de
centrale onderzoeksvraag met betrekking tot radiogebaseerde bestandstoegang te 
beantwoorden. Om de technische haalbaarheid en de operationele grenzen van een 
IP-over-radioverbinding met een private fileserver objectief te kunnen kwantificeren, is 
het onderzoek opgedeeld in vier opeenvolgende, logisch met elkaar verbonden fasen. 
Elke fase kent een specifieke doelstelling, onderzoeksdoel en tastbare deliverable. 
Dit gestructureerde proces waarborgt dat de uiteindelijke empirische metingen direct 
kunnen worden getoetst aan vooraf gedefinieerde, harde netwerkeisen.

\section{Fase 1: Literatuurstudie (State-of-the-Art)}
\label{sec:meth-fase1}

De eerste fase had als doel het creëren van een sluitend theoretisch en conceptueel kader.
Radiogebaseerde datacommunicatie bevindt zich op het snijvlak van informatietheorie en 
toegepaste netwerkarchitectuur, waardoor een diepgaande theoretische verkenning van 
de fysieke laag en de transportlaag noodzakelijk was.

\begin{itemize}
    \item \textbf{Onderzoeksmethode:} Systematisch literatuuronderzoek (*desk research*) 
    waarbij gebruik is gemaakt van peer-reviewed artikelen, academische handboeken en 
    officiële specificatiedocumenten van instanties zoals de ITU en het BIPT.
    \item \textbf{Inhoudelijke focus:} Analyse van de wet van Shannon-Hartley, atmosferische 
    propagatieverschijnselen (troposferische ducting en Sporadic-E), robuuste modulatieschema's 
    (OFDM versus single-carrier QPSK), netwerkprotocollen (AX.25, Winlink, native IP-over-radio) 
    en de fysieke limieten van specifieke spectrums (zoals de 60~GHz zuurstofabsorptie).
    \item \textbf{Deliverable:} Het reeds uitgewerkte Hoofdstuk~\ref{ch:stand-van-zaken} 
    (Stand van zaken), dat de theoretische randvoorwaarden voor de rest van dit onderzoek vastlegt.
\end{itemize}

\section{Fase 2: Requirementsanalyse en Spectrumvergelijking}
\label{sec:meth-fase2}

In de tweede fase is de transitie gemaakt van abstracte theorie naar concrete, meetbare 
systeemeisen. Voordat een fysieke testopstelling kan worden ontworpen, moeten de 
minimale prestatie-indicatoren van het netwerk en de meest geschikte frequentiebanden 
worden gedefinieerd.

\begin{itemize}
    \item \textbf{Onderzoeksmethode:} Requirementsanalyse gecombineerd met een 
    kwalitatieve vergelijkende studie. De eisen zijn opgesteld volgens de industriestandaard 
    voor netwerkprotocollen (RFC-stijl).
    \item \textbf{Inhoudelijke focus:} Het definiëren van harde netwerktoleranties 
    met betrekking tot minimale doorvoersnelheid (*throughput*), maximaal latency- en jitterbudget, 
    en het maximaal toelaatbare packet loss-percentage waaronder stateful transportprotocollen 
    (zoals TCP) stabiel blijven functioneren. Tegelijkertijd worden de functionele trade-offs 
    tussen gelicentieerde radioamateurfrequenties en licentievrije ISM-banden gecatalogiseerd.
    \item \textbf{Deliverable:} Een gestructureerde architectuur- en requirementsmatrix die de 
    basis vormt voor het hierop volgende Hoofdstuk~\ref{ch:requirementsanalyse} (Requirementsanalyse).
\end{itemize}

\section{Fase 3: Ontwerp en Realisatie van de Proof-of-Concept}
\label{sec:meth-fase3}

De derde fase beslaat het praktische en experimentele kernonderdeel van deze bachelorproef. 
Om te verifiëren of bestandsoverdrachtsprotocollen daadwerkelijk kunnen functioneren over 
een draadloos radiomedium, is een representatieve fysieke testomgeving onmisbaar.

\begin{itemize}
    \item \textbf{Onderzoeksmethode:} Experimenteel-technisch ontwerponderzoek (*design science research*). 
    Er is bewust gekozen voor het bouwen van een fysieke proof-of-concept (PoC) in plaats van 
    een puur gesimuleerde netwerkomgeving, aangezien reële atmosferische demping en 
    hardware-overhead via simulatiesoftware moeilijk betrouwbaar te modelleren zijn.
    \item \textbf{Inhoudelijke focus:} Selectie en configuratie van de radiohardware (bijvoorbeeld 
    radiomodems of gerichte Point-to-Point straalantennes) en de netwerkinfrastructuur. Dit omvat 
    het opzetten van een transparante IP-tunnel over de radioverbinding, het inrichten van een 
    Linux-gebaseerde fileserver (met protocollen zoals SFTP, SMB of WebDAV) en het configureren 
    van de clientnode.
    \item \textbf{Deliverable:} Een functionele, infrastructuuronafhankelijke netwerkopstelling, 
    volledig gedocumenteerd in Hoofdstuk~\ref{ch:proof-of-concept} (Ontwerp en Realisatie).
\end{itemize}

\section{Fase 4: Empirische Evaluatie en Data-analyse}
\label{sec:meth-fase4}

In de finale fase zijn gecontroleerde praktijktests uitgevoerd om de prestaties van de 
gebouwde proof-of-concept kwantitatief in kaart te brengen en te valideren.

\begin{itemize}
    \item \textbf{Onderzoeksmethode:} Kwantitatief empirisch experimenteel onderzoek. 
    Door middel van systematische stresstests zijn de netwerkprestaties gemeten onder 
    wisselende onafhankelijke variabelen (zoals afstand, bestandstype en transmissie-omstandigheden).
    \item \textbf{Inhoudelijke focus:} Het verzamelen van harde netwerkmetrieken met behulp van 
    diagnostische netwerktools. De primaire focus ligt op het meten van de effectieve bitsnelheid 
    bij variërende bestandsgroottes, de invloed van packet loss op de TCP-venstergrootte (*window size*), 
    en de totale overdrachttijd. Hierbij worden tevens praktijkobservaties omtrent langere-afstandsverbindingen 
    (zoals troposferische propagatie-invloeden op de UHF-band) geanalyseerd.
    \item \textbf{Deliverable:} Een statistische evaluatie, grafische prestatieplots en een finaal 
    haalbaarheidsoordeel, uitgewerkt in Hoofdstuk~\ref{ch:resultaten-en-evaluatie} (Resultaten en Evaluatie).
\end{itemize}

\section{Verantwoording van de Aanpak}
\label{sec:meth-verantwoording}

De keuze voor een opeenvolgende, fasegewijze benadering is essentieel om de interne 
en externe validiteit van dit toegepaste IT-onderzoek te garanderen. Radiokanalen zijn van nature 
stochastisch en onderhevig aan omgevingsfactoren, wat een strikte isolatie van parameters 
vereist. Door eerst de theoretische limieten te analyseren (Fase 1) en om te zetten in harde 
netwerkeisen (Fase 2), wordt voorkomen dat de daaropvolgende testfase (Fase 3) ad-hoc of 
onvolledig verloopt. De empirische data-analyse (Fase 4) fungeert tot slot als de ultieme, 
objectieve verificatie. Deze methodologische keten zorgt ervoor dat de uiteindelijke 
conclusies van deze bachelorproef niet berusten op anekdotische praktijkervaringen, maar op 
reproduceerbare, wetenschappelijk onderbouwde meetresultaten.

